\documentclass[11pt]{article}
\usepackage[margin=1in]{geometry}
\usepackage{amsmath,amssymb,amsthm,mathtools}
\usepackage{graphicx}
\usepackage{microtype}
\usepackage[hidelinks]{hyperref}

\newtheorem{theorem}{Theorem}
\newtheorem{lemma}{Lemma}
\newcommand{\LP}{\mathcal{L}\!\mathcal{P}}

\title{A Paired-Zero Subcase of the Br\"anden--Chasse\\Fourier Transform Conjecture}
\author{Candidate partial result for arXiv:1402.2795}
\date{27 August 2026}

\begin{document}
\maketitle

\begin{center}
\fbox{\parbox{0.9\linewidth}{\textbf{Status.} Candidate partial result,
likely valid.  The unrestricted conjecture is not claimed.}}
\end{center}

\begin{abstract}
Br\"anden and Chasse conjecture that if a real transcendental entire function
$G$ has derivative in the Laguerre--P\'olya class and decays appropriately on
the imaginary axis, then the Fourier transform of $e^{G(it)}$ is again in the
Laguerre--P\'olya class.  We prove the conjecture when the zeros of $G'$ admit a
summable positive/negative pairing whose total imaginary-axis angular defect
is less than $\pi/2$.  The proof constructs real-rooted polynomial derivative
approximants and obtains a uniform Gaussian-or-better tail estimate, which is
the compactness input missing in the unrestricted problem.
\end{abstract}

\section{Source question}

The target is Conjecture 6.1 on page 29 of Br\"anden--Chasse
\cite{BrandenChasse}.  The source statement is reproduced below.

\begin{center}
\includegraphics[width=0.96\linewidth]{figures/open_problem_crop.png}
\end{center}

The displayed source uses the notation $G(it)\to-\infty$.  Since complex
numbers are not ordered, and since Corollary 5.5 and Proposition 5.6 directly
before the conjecture use $\operatorname{Re}G(it)\to-\infty$, we treat the
nontrivial intended reading as
\[
  \operatorname{Re}G(it)\longrightarrow-\infty
  \qquad(t\to\pm\infty).
\]

Recall that $\LP$ is the class of real entire functions that are locally
uniform limits of real polynomials having only real zeros.

\section{Candidate partial theorem}

\begin{theorem}[paired genus-zero subcase]\label{thm:main}
Let $m\in\mathbb{Z}_{\ge0}$, $\alpha\ge0$, and $c\in\mathbb{R}$ satisfy
$c(-1)^m>0$.  Let $a_j,b_j>0$ satisfy
\begin{equation}\label{eq:summability}
 \sum_{j=1}^{\infty}\left(\frac1{a_j}+\frac1{b_j}\right)<\infty
\end{equation}
and
\begin{equation}\label{eq:angle}
 A:=\sum_{j=1}^{\infty}
 \arctan\!\left(\frac{|a_j-b_j|}{2\sqrt{a_jb_j}}\right)<\frac{\pi}{2}.
\end{equation}
Define
\begin{equation}\label{eq:H}
 H(z)=c z^{2m+1}e^{-\alpha z^2}
 \prod_{j=1}^{\infty}\left(1-\frac{z}{a_j}\right)
                         \left(1+\frac{z}{b_j}\right),
\end{equation}
and let $G$ be the real entire primitive of $H$ with a prescribed real value
$G(0)$.  Then
\begin{equation}\label{eq:coercive}
 \operatorname{Re}G(it)
 \le G(0)-\frac{|c|\cos A}{2m+2}|t|^{2m+2}
 \qquad(t\in\mathbb{R}).
\end{equation}
Consequently
\[
 F(z)=\int_{-\infty}^{\infty}e^{G(it)}e^{izt}\,dt
\]
is an entire function in $\LP$.
\end{theorem}

Condition \eqref{eq:summability} makes \eqref{eq:H} a convergent genus-zero
canonical product.  In particular $H\in\LP$.  The conclusion therefore proves
the source conjecture for an infinite-dimensional, explicitly checkable
paired-zero class.

\section{Idea of the proof}

For one paired positive/negative zero, substitution $z=it$ produces a factor
with positive real part.  Its argument can deviate from zero by at most
$\arctan(|a_j-b_j|/(2\sqrt{a_jb_j}))$.  The strict total bound
\eqref{eq:angle} therefore keeps every finite canonical-product truncation of
$H(it)$ inside the same open half-plane.  Its imaginary part has the sign of
$t$ and a uniform polynomial lower bound.

After integration, this becomes one coercive upper bound for the real parts of
all polynomial primitives.  Br\"anden--Chasse's polynomial theorem applies to
each primitive.  The common coercive bound is strong enough for dominated
convergence of the Fourier transforms, so the Laguerre--P\'olya conclusion
survives the transcendental limit.

\section{Proof}

For $N\ge1$, put
\begin{equation}\label{eq:HN}
 H_N(z)=c z^{2m+1}\left(1-\frac{\alpha z^2}{N}\right)^N
 \prod_{j=1}^{N}\left(1-\frac{z}{a_j}\right)
                    \left(1+\frac{z}{b_j}\right).
\end{equation}
(When $\alpha=0$, the middle factor is $1$.)  Every $H_N$ is a real
polynomial with only real zeros, and $H_N\to H$ locally uniformly.  Let $G_N$
be its primitive normalized by $G_N(0)=G(0)$.  Then $G_N\to G$ locally
uniformly.

Fix $t>0$.  The $j$th paired factor in \eqref{eq:HN} equals
\begin{align*}
 \left(1-\frac{it}{a_j}\right)\left(1+\frac{it}{b_j}\right)
 &=1+\frac{t^2}{a_jb_j}
   +it\left(\frac1{b_j}-\frac1{a_j}\right).
\end{align*}
Its real part is positive.  If $\theta_j(t)$ is its principal argument, then
\[
 |\tan\theta_j(t)|
 =\frac{t|a_j-b_j|}{a_jb_j+t^2}
 \le\frac{|a_j-b_j|}{2\sqrt{a_jb_j}},
\]
because $t/(a_jb_j+t^2)$ is maximized at $t=\sqrt{a_jb_j}$.  Hence
\begin{equation}\label{eq:theta}
 |\theta_j(t)|\le
 \arctan\!\left(\frac{|a_j-b_j|}{2\sqrt{a_jb_j}}\right).
\end{equation}

The factor $(1+\alpha t^2/N)^N$ is positive.  Moreover
$c(it)^{2m+1}=i|c|t^{2m+1}$ because $c(-1)^m>0$.  Equations
\eqref{eq:angle} and \eqref{eq:theta} therefore give
\begin{equation}\label{eq:imagbound}
 \operatorname{Im}H_N(it)
 \ge |c|\cos(A)t^{2m+1},\qquad t>0.
\end{equation}
Indeed, the argument of $H_N(it)$ differs from $\pi/2$ by at most $A$, while
every remaining factor in its modulus is at least $1$.  Since $H_N$ has real
coefficients, conjugation gives the corresponding negative sign for $t<0$.

Along the imaginary axis,
\[
 \frac{d}{dt}\operatorname{Re}G_N(it)
 =\operatorname{Re}\bigl(iH_N(it)\bigr)
 =-\operatorname{Im}H_N(it).
\]
Integrating \eqref{eq:imagbound} on either side of the origin yields the
uniform estimate
\begin{equation}\label{eq:GNbound}
 \operatorname{Re}G_N(it)
 \le G(0)-q|t|^{2m+2},
 \qquad q:=\frac{|c|\cos A}{2m+2}>0.
\end{equation}
Passing to the pointwise limit gives \eqref{eq:coercive} for $G$.

Each $G_N$ is a real polynomial, $G_N'=H_N\in\LP$, and
$\operatorname{Re}G_N(it)\to-\infty$ at both ends.  Corollary 5.5 of
\cite{BrandenChasse} therefore implies
\[
 F_N(z):=\int_{-\infty}^{\infty}e^{G_N(it)}e^{izt}\,dt\in\LP.
\]
If $K\subset\mathbb{C}$ is compact and
$M=\sup_{z\in K}|\operatorname{Im}z|$, then \eqref{eq:GNbound} gives the
$N$-independent integrable majorant
\[
 |e^{G_N(it)}e^{izt}|
 \le \exp\bigl(G(0)-q|t|^{2m+2}+M|t|\bigr),
 \qquad z\in K.
\]
Dominated convergence now yields $F_N\to F$ locally uniformly on
$\mathbb{C}$.  The class $\LP$ is closed under locally uniform limits.  The
limit is not identically zero, since $e^{G(it)}$ is a nonzero integrable
function and the Fourier transform is injective on $L^1(\mathbb{R})$.
Thus $F\in\LP$, completing the proof. \qed

\section{Verification notes and limitations}

The proof uses no numerical claim.  The phase estimate, the sign for negative
$t$, and the primitive derivative were recalculated independently; details are
recorded in the companion verification report.

The added hypothesis \eqref{eq:angle} is the essential limitation.  A general
Laguerre--P\'olya function has locally convergent real-rooted polynomial
approximants, but their primitives need not share an integrable imaginary-axis
majorant.  The source condition controls the accumulated primitive, not the
pointwise phase of arbitrary truncations.  Thus this packet does not resolve
the unrestricted conjecture.

A bounded search of the run indexes, exact web/arXiv phrases, and the indexed
citing works of the published source found no later proof of Conjecture 6.1
and no occurrence of criterion \eqref{eq:angle}.  This is not a definitive
novelty claim.  Some functions covered here may also satisfy Proposition 5.6
of the source; strict separation from that coefficient-sign class has not been
established.

\paragraph{Human-review recommendation.}
Check first the intended $\operatorname{Re}G(it)$ interpretation and the use of
source Corollary 5.5.  Conditional on that textual interpretation, the partial
theorem itself should be reviewed as a short canonical-product and dominated-
convergence argument.

\begin{thebibliography}{9}
\bibitem{BrandenChasse}
P.~Br\"anden and M.~Chasse,
\emph{Classification theorems for operators preserving zeros in a strip},
J. Anal. Math. \textbf{132} (2017), 177--215;
arXiv:1402.2795.
\end{thebibliography}

\end{document}
